% Owned by #24 (dataset and methodology). The split table is generated from
% data/splits.csv by scripts/build_tables.py --table dataset.

\subsection{Dataset}
\label{sec:dataset}

We use EuroSAT~\cite{helber2019eurosat}, a land-use and land-cover dataset
containing 27,000 labelled, geo-referenced Sentinel-2 image patches across ten
classes: AnnualCrop, Forest, HerbaceousVegetation, Highway, Industrial,
Pasture, PermanentCrop, Residential, River, and SeaLake. Although the original
dataset includes all 13 Sentinel-2 spectral bands, this work uses the RGB
version, which contains the three visible bands stored as 8-bit images. The
dataset is available under the MIT licence and was downloaded through
\texttt{torchvision}~\cite{torchvision2016} version 0.20 or later, which uses
the maintained Hugging Face mirror. Because each class represents a land-cover
category viewed from an overhead perspective, orientation changes such as
flips and 90-degree rotations do not alter the class label, which supports the
augmentation choices described in Section~\ref{sec:data-preparation}.

\paragraph{Input resolution.}
All EuroSAT images have a native resolution of $64\times64$ pixels. Since the
assignment specifies a maximum input resolution of $64\times64$ for the
low-resource setting, the dataset already satisfies this requirement and no
downscaling is needed. The custom models therefore use the original images
directly, without reducing their spatial resolution.
Figure~\ref{fig:data-augmentation-grid} shows one native sample from each
class. The only resizing beyond the native resolution is an upsampling step to
$224\times224$ for the supplementary SOTA resolution experiment in Task~5,
where the pretrained models are also evaluated at their standard design
resolution.

\paragraph{Split.}
The dataset was divided once into training, validation, and test sets using
the required 70/15/15 proportions. The split was stratified by class, so each
class was divided independently and the class distribution remains consistent
across all three sets. Class sizes range from 2,000 to 3,000 images
(Table~\ref{tab:dataset-split}), and these sizes allow the 70/15/15
proportions to be applied exactly. Within each class, image files were sorted
by name and then shuffled using a fixed seed of 42. This makes the split
reproducible and ensures that the same partition is obtained across different
machines.

\paragraph{One partition for every experiment.}
The train, validation, and test assignment for each image is recorded in the
committed manifest, \texttt{data/splits.csv}. All experiments use this same
file, including the custom models, optimizer comparison, and fine-tuned SOTA
models, rather than generating new splits. This keeps the data partition fixed
across experiments and satisfies the requirement that the SOTA models use the
same splits as the custom networks. Table~\ref{tab:dataset-split} is generated
directly from this manifest. The validation set is used for model and
optimizer selection, while the test set is reserved for final evaluation after
training.

\begin{table}[htbp]
  \centering
  \caption{Images per class in each set, generated from the committed split
  manifest \texttt{data/splits.csv}.}
  \label{tab:dataset-split}
  \generatedtable{tables/dataset.tex}{\#24}
\end{table}
